%chapter11
\chapter[常用教程]{常用教程}

\section[校园网]{校园网}
\begin{enumerate}
    \item 激活：到校报到完成后，学校将按照个人身份证号在\uline{\href{http://210.44.80.65/}{http://210.44.80.65/}}为大家开通校园网账号，并告知初始密码\footnotemark
          \footnotetext{每年不同，详见班级群内具体通知。}
    \item 充值：点击\uline{\href{https://slzfw.wfmc.edu.cn}{https://slzfw.wfmc.edu.cn}}，并登录；或在校园网的登陆页面点击“自助服务”按钮\footnotemark，登陆后即可使用支付宝充值（微信支付暂时无法使用）
          \footnotetext{若点击按钮后网页显示“403 Error”，请将网址最前面的“http”改为“https”，按回车即可。}
    \item 校园网流量当前政策为每月免费60G，超出部分$0.5￥/G$
\end{enumerate}

\section[校园手机卡]{校园手机卡}
\begin{enumerate}
    \item \textbf{开通与否全凭自愿}，随录取通知书一并寄出
    \item 通常内含200分钟全国通话、10G全国流量、80G校园流量（仅山东省内所有高校可用）
    \item 开学报道当天可前往大服，免费领取礼品
\end{enumerate}

\section[空调使用教程]{空调使用教程}
\begin{enumerate}
    \item 微信关注“海享租”公众号，点击公众号菜单“在线租赁”，并注册、登录
    \item 点击“扫一扫”功能，扫描空调右下角二维码进行租赁
    \item 租赁完成后，点击“设备”→“空调图标”→“时长”，进行充值
    \item 点击“设备”→“空调图标”→“成员管理”，在此功能下设置宿舍全部成员均可管理空调开关
    \item 注意：空调使用时长收费（0.55元/小时），具体收费政策详见“海享租”公众号
\end{enumerate}

\section[浴室预约与使用]{浴室预约与使用}
\begin{enumerate}
    \item 安装：
          \begin{enumerate}
              \item 在手机应用市场下载“大白U帮”app
              \item 按照实际住宿情况注册
              \item 授予并开启“定位”与“蓝牙”权限
          \end{enumerate}
    \item 本楼层小浴室使用：
          \begin{enumerate}
              \item 带好洗浴物品前往公共厕所旁边的浴室排队
              \item 进入浴室，点击如右图所示的按钮（\mbox{\includegraphics[height=2.4ex]{bath.png}}）
              \item 选择“蓝牙设备”→“点击进行时”→“洗澡”
              \item →“搜索洗澡”\footnotemark →选择设备\footnotemark →“开始洗澡”
              \footnotetext[3]{搜索不到设备请务必开启“蓝牙”功能，学校的设备无法扫码连接。}
              \footnotetext{距离入口最近的是1号，远的是2号；不确定可以询问学长。}
              \item 结束后点击“结束洗澡”按钮，并结算
          \end{enumerate}
    \item 一层公共大浴室预约：
          \begin{enumerate}
              \item 在软件初始界面根据实际情况选择“X号楼1层”的浴室
              \item 点击一个浴位，并点击“预约”按钮（若已满请选择“排队”）
              \item 按照软件提示前往洗浴
              \item 结束后点击“停止”按钮，并结算
          \end{enumerate}
    \item 费用：以程序显示为准，详情收费标准略
    \item 申诉：如果因洗澡时突然停电导致无法结束洗澡而被扣费，请按照软件打开时弹出的公告，联系相关工作人员处理
    \item \textbf{注意：如果未点击“结束洗澡”的按钮可能会被多扣费}
\end{enumerate}

\section[洗衣机使用教程]{洗衣机使用教程}
\begin{enumerate}
    \item 前往宿舍一层，位于公共厕所旁的洗衣间，微信扫码注册后使用即可
    \item 需要自带洗衣液（极个别同学可能会使用洗衣机洗鞋）
    \item 收费标准详见软件提示
\end{enumerate}

\section[吹风机使用教程]{吹风机使用教程}
\begin{enumerate}
    \item 每层公共浴室旁边有两个公用吹风机，需扫码\footnotemark 租赁使用
    \footnotetext{部分吹风机屏幕二维码可能有缺损，不易扫描成功，多次尝试即可。}
    \item 待手机发出“滴”-“滴”的声音后租赁成功
    \item 将手机扬声器对准吹风机租赁器方可正常使用
    \item 收费标准：详见软件提示，1分钱起步
\end{enumerate}

\section[设施报修方式]{设施报修方式}
\begin{enumerate}
    \item 加入各宿舍楼的QQ报修群，在群内反映具体故障
    \item 前往一层宿管处填表报修
    \item 在宿舍一层宿管旁边的公告栏处查看负责该楼层的物业电话，直接拨打即可
\end{enumerate}

\section[钉钉请假流程]{钉钉请假流程\footnotemark}
\footnotetext{待学校统一将大家拉入钉钉的“潍坊医学院”企业后方可使用。}
\begin{enumerate}
    \item 注意：各学院要求不一，仅以临床医学院为例
    \item 线下请假步骤（正常情况）：
    \begin{enumerate}
        \item 前往学工办或班主任办公室，当面请假并获得假条
        \item 根据老师要求扫描相关二维码
        \item 钉钉填表
        \item 扫脸进出校门，并将请假条之一交给保卫处
        \item 返校后，在钉钉的电子假条处，以评论的方式销假
    \end{enumerate}
    \item 线上请假步骤：
    \begin{enumerate}
        \item 打开钉钉→点击左上角选择主企业为“潍坊医学院”
        \item →点击页面最下方菜单栏“工作台”→点击“OA审批”
        \item →在“学生日常事务管理”类选择“浮烟山校区本科学生请假单”→填表
        \item 电话联系班主任老师或学工办老师，说明请假事由并等待审批
        \item 审批通过后即可扫脸进出校门
        \item 返校后，在钉钉的电子假条处，以评论的方式销假
    \end{enumerate}
\end{enumerate}

\section[学工系统（微信小程序）]{学工系统（微信小程序）}
\begin{enumerate}
    \item 用途：学工系统主要用于晚点名、返校信息填报等日常工作\footnotemark
    \footnotetext{原“请假审批、外出审批”工作已基本转移至“钉钉”，详见\uline{text}。}
    \item 使用方式：
    \begin{enumerate}
        \item 在微信搜索“智慧学工”小程序
        \item 授予“定位”权限
        \item 根据学校下发的账号密码进行登录（推荐立即与微信绑定以免忘记密码）
    \end{enumerate}
\end{enumerate}

\section[教务系统]{教务系统}
\begin{enumerate}
    \item 官网：\uline{\href{https://jwgl.wfmc.edu.cn/}{https://jwgl.wfmc.edu.cn/}}
    \item 用途：\textbf{选课，缓考申请，成绩查询}，查看（导出）课程表
    \item 缓考申请教程：
    \begin{enumerate}
        \item 进入教务系统
    \end{enumerate}
    \item 注意：仅限校内访问，如需在家使用教务系统，请参见\uline{\ref{cas_system}}条目
\end{enumerate}

\section[资源访问控制系统]{资源访问控制系统}
\label{cas_system}

1.本系统用于在校外访问校内网络信息资源，如：知网、教务系统等；
2.登录 https://webvpn.wfmc.edu.cn/login ，点击“统一身份认证登录”按钮，按照学校下发的专用账号、密码登录即可；
3.如需在校外访问教务系统，请完成上述操作后选择“教务系统-非单点登录”；
4.如需校外查阅文献推荐使用CARSI系统，账密同校园网账密；
5.因本系统响应速度较慢，推荐在校内直接登录教务系统、知网等的官网进行操作。

（十一）校务行（微信小程序）
1.用于快速成绩查询、成绩报告单打印、学籍证明打印；
2.打印的收费标准详见小程序内部说明。

（十二）常用课程表程序
1.wakeup课程表：各手机应用商店直接下载，按教程进行即可，无广告；
2.超级课程表：各手机应用商店直接下载，有广告，使用人数较多。

（十三）学生票购买流程（具体规定详见12306官网）
1.入学后，根据学校统一通知，网上填报“家乡⇄学校”的最近的火车站点信息；
2.学生证注册章盖章后（每学期均需要），拿学生证、身份证前往任一火车站自动售票机处，选择“学生优惠资质绑定”按钮，按提示进行操作；或直接在人工售票处完成学生优惠资质认证（注意：每年均需重新认证学生资质一次）；
3.全年有4次优惠机会，高铁75折，普通火车5折。

（十四）文体中心预约教程
1.关注“潍医文体中心”公众号；
2.使用cas认证系统的账号密码登录，并录入面部信息，预约并缴费后即可前往；
3.收费情况详见公众号文章，校外人员可在有余票时直接千万前往前台购票入场。